% Chapter 10 of The Replica Rosetta Stone: the LASSO and robust regression.
\rstchapter{The LASSO and Robust Regression}{Four routes to one fixed point}{the replica fixed point\dict AMP and the CGMT}{ch:lasso}
\stonefoot{proved}
\begin{refsection}[chapters/10-lasso.bib]

\begin{chaptersummary}
When a sparse vector is estimated from Gaussian linear measurements whose number is proportional to the dimension, the LASSO behaves, coordinate by coordinate, like soft thresholding of the truth seen in Gaussian noise, with a noise level and a threshold fixed by two scalar equations. Physicists derived this fixed point with zero-temperature replicas; Bayati and Montanari proved it with approximate message passing, and Gordon's comparison proves it again, through an auxiliary problem whose scalar form is the replica-symmetric potential after a change of variables. We prove the limit of the LASSO's error by the convex Gaussian min-max theorem and the noiseless Donoho--Tanner threshold by computing a statistical dimension, and we show that the replica, state-evolution and geometric formulas for that threshold are one function. Simulations of the LASSO, of AMP, of $\ell_1$ minimization and of Huber regression agree with every prediction; the open questions concern non-Gaussian designs, non-convex penalties and finite-size corrections.
\end{chaptersummary}

\section{The question: what does sparsity buy?}\label{sec:lasso:question}

A statistician with $n$ observations and $d>n$ candidate predictors cannot fit least squares: the minimum-norm interpolant of \cref{ch:ridge} keeps only the fraction $n/d$ of the signal. If few coefficients are nonzero, the \term{LASSO}\index{LASSO} of \textcite{Tibshirani1996} adds an $\ell_1$ penalty and sets many coordinates exactly to zero,\sidenote{Tibshirani wrote the penalty as a constraint, $\norm\beta_1\le t$; by convex duality the constrained and penalized forms trace the same path of solutions.} and its noiseless limit, $\ell_1$ minimization or basis pursuit \parencite{ChenDonohoSaunders1998}, can recover a sparse vector from fewer measurements than unknowns. How many measurements are needed, and how large is the error when they are noisy, if $n$ and $d$ are large and comparable?

We use the scaling of \textcite{BayatiMontanari2012}, in which the unknown has entries of order one. The design $A\in\R^{n\times d}$ has independent $\Normal(0,1/n)$ entries, so $A=X/\sqrt n$ for the design $X$ of \cref{ch:ridge}; the unknown is $\theta_0=\sqrt n\,\beta^*$; the noise $w$ has independent $\Normal(0,\sigma^2)$ entries (other chapters write it $\sigma\varepsilon$; here $\varepsilon$ is a sparsity); $y=A\theta_0+w$ and $\alpha=n/d$. The LASSO is
\begin{equation}\label{eq:lasso:lasso}
\hat\theta\in\argmin_{\theta\in\R^d}F(\theta),\qquad F(\theta)=\tfrac12\norm{y-A\theta}^2+\lambda\norm\theta_1,\quad\lambda>0,
\end{equation}
and its error is $\mathrm{MSE}=\norm{\hat\theta-\theta_0}^2/d$. Since $\hat\theta=\sqrt n\,\hat\beta$, the excess risk of the other regression chapters is $R=\norm{\hat\beta-\beta^*}^2=\mathrm{MSE}/\alpha$; this is the only translation the chapter needs.\sidenote{In the per-sample normalization $\frac1{2n}\norm{y-X\beta}^2+\lambda'\norm\beta_1$ of \cref{ch:ridge} the same estimator has $\lambda'=\lambda/\sqrt n$, which tends to zero; hence the scaling used here.} The entries of $\theta_0$ are independent with law $P_{\Theta_0}$, and $\varepsilon=\Prob(\Theta_0\ne0)$. The running example is the Bernoulli--Gaussian signal $(1-\varepsilon)\delta_0+\varepsilon\Normal(0,1)$ with $\varepsilon=0.1$, at $\alpha=0.5$ and $\sigma^2=0.05$.

\term{Soft thresholding}\index{soft thresholding}, $\eta(u;\zeta)=\sgn(u)(\abs u-\zeta)_+$, minimizes $\frac12(x-u)^2+\zeta\abs x$: for $x\ne0$ the condition $x-u+\zeta\sgn x=0$ has a solution exactly when $\abs u>\zeta$. If the columns of $A$ were orthonormal, which needs $n\ge d$, the LASSO would be $\eta(A^\top y;\lambda)$ coordinatewise, and $A^\top y=\theta_0+A^\top w$ is the truth plus Gaussian noise of variance $\sigma^2$. The physical prediction is that this picture survives at $n<d$, once the noise and the threshold are renormalized.

\begin{physicsbox}[replicas and cavities, 2009--2016]
Let $n,d\to\infty$ with $n/d\to\alpha$. \textup{(i)} The empirical distribution of the pairs $(\hat\theta_j,\theta_{0j})$ tends to the law of $(\eta(\Theta_0+\tau Z;\zeta),\Theta_0)$, with $Z\sim\Normal(0,1)$ independent of $\Theta_0$, where $\tau>0$ and $\zeta>0$ solve
\begin{equation}\label{eq:lasso:fp}
\tau^2=\sigma^2+\frac1\alpha\E\bigl[(\eta(\Theta_0+\tau Z;\zeta)-\Theta_0)^2\bigr],\qquad
\lambda=\zeta\Bigl(1-\frac{\Prob(\abs{\Theta_0+\tau Z}>\zeta)}\alpha\Bigr).
\end{equation}
So $\mathrm{MSE}\to\alpha(\tau^2-\sigma^2)$, and the fraction of nonzero coordinates of $\hat\theta$ tends to $\Prob(\abs{\Theta_0+\tau Z}>\zeta)$. \textup{(ii)} Without noise, $\ell_1$ minimization recovers $\theta_0$ when $\alpha>\psi(\varepsilon)$ and fails when $\alpha<\psi(\varepsilon)$, where $D(t)=\E(\abs Z-t)_+^2$ and
\begin{equation}\label{eq:lasso:psi}
\psi(\varepsilon)=\min_{t\ge0}J_\varepsilon(t),\qquad J_\varepsilon(t)=\varepsilon(1+t^2)+(1-\varepsilon)D(t).
\end{equation}
\textup{(iii)} The M-estimator $\argmin_\theta\sum_i\rho(y_i-\ip{a_i}\theta)$, with a convex loss $\rho$, the rows $a_i$ of $A$, noise $w_i\iid P_W$ and $\alpha>1$, has $R\to r^2$, where, with $\hat z=W+rZ$, $W\sim P_W$ independent of $Z$, and $\operatorname{prox}_c(x)=\argmin_u\{c\rho(u)+\frac12(u-x)^2\}$,
\begin{equation}\label{eq:lasso:mest}
\E\bigl[\operatorname{prox}_c'(\hat z)\bigr]=1-\frac1\alpha,\qquad \frac{r^2}\alpha=\E\bigl[(\hat z-\operatorname{prox}_c(\hat z))^2\bigr].
\end{equation}
\end{physicsbox}

Prediction (i), the \term{decoupling principle}\index{decoupling principle}, is the replica-symmetric result of \textcite{RanganFletcherGoyal2012}; \eqref{eq:lasso:fp} is the fixed point of \textcite{BayatiMontanari2012}. Prediction (ii) was derived by replicas \parencite{KabashimaEtAl2009} and by state evolution \parencite{DonohoMalekiMontanari2009}, and earlier, as a bound, from the faces of random polytopes \parencite{Donoho2006,DonohoTanner2009a}; the curve is the \term{Donoho--Tanner phase transition}\index{Donoho--Tanner phase transition}, plotted in the coordinates $(\delta,\rho)=(\alpha,\varepsilon/\alpha)$. Prediction (iii) concerns the \term{M-estimators}\index{M-estimator} of robust regression \parencite{Huber1964}; \textcite{Huber1973} already let the number of parameters grow with $n$, and the system was derived by a probabilistic heuristic \parencite{ElKarouiEtAl2013} and by replicas \parencite{AdvaniGanguli2016}. It says that the estimator sees the noise inflated by an independent Gaussian of variance $r^2$; for least squares it gives $R=\sigma^2/(\alpha-1)$, which a robust loss can beat under heavy tails.

How do physicists derive the fixed point, and what is proved? Every item of the box is a theorem for Gaussian designs; we prove the error of the LASSO and the noiseless threshold in full. Four routes reach the same fixed point: replicas, the state evolution of approximate message passing, Gordon's comparison and, without noise, the geometry of convex cones.

\section{The physics: replicas at zero temperature}\label{sec:lasso:physics}

The derivation is the standard one, in this book's normalization: \textcite{KabashimaEtAl2009} gave it for noiseless $\ell_p$ reconstruction (their equation~(7) is the noiseless form of \eqref{eq:lasso:rs} below), \textcite{RanganFletcherGoyal2012} with noise, \textcite{AdvaniGanguli2016} for general convex losses and penalties. The replica trick is introduced in \cref{ch:freeentropy}, and its zero-temperature version for ridge in \cref{ch:ridge}.

\begin{heuristic}[The replicated LASSO]
\emph{Step 1: the design.} Let $Z_\beta=\int e^{-\beta F(\theta)}\dd\theta$; as $\beta\to\infty$ the Gibbs measure $e^{-\beta F}/Z_\beta$ concentrates on the minimizers, and $e_0=\lim_d\frac1d\min F$ is the limit of $-\frac1{\beta d}\E\log Z_\beta$, computed from the moments $\E Z_\beta^k$ continued to $k\to0$. With replicas $\theta^1,\dots,\theta^k$ and $v^a=\theta^a-\theta_0$, the identity $y-A\theta^a=w-Av^a$ shows that the design enters only through the fields $u_i^a=\ip{a_i}{v^a}$, independent over $i$ and centered Gaussian with covariance $\ip{v^a}{v^b}/n=(q_{ab}-m_a-m_b+\E\Theta_0^2)/\alpha$, where $q_{ab}=\ip{\theta^a}{\theta^b}/d$ and $m_a=\ip{\theta^a}{\theta_0}/d$ are the overlaps.

\emph{Step 2: replica symmetry.} Assume $q_{aa}=Q$, $q_{ab}=q$ ($a\ne b$) and $m_a=m$. Then $u_i^a=\sqrt\omega\,\xi_0+\sqrt{(Q-q)/\alpha}\,\xi_a$ with independent standard Gaussians $\xi_a$ and $\omega=(q-2m+\E\Theta_0^2)/\alpha$, and with $s=w_i-\sqrt\omega\,\xi_0\sim\Normal(0,\sigma^2+\omega)$ and $\Delta_\beta=\beta(Q-q)/\alpha$ a Gaussian integral gives the rows' contribution, per coordinate,
\begin{multline*}
\alpha\log\E_s\Bigl[\bigl(\E_\xi e^{-\beta(s-\sqrt{(Q-q)/\alpha}\,\xi)^2/2}\bigr)^k\Bigr]\\
=-\frac{\alpha k}2\log(1+\Delta_\beta)-\frac\alpha2\log\Bigl(1+\frac{k\beta(\sigma^2+\omega)}{1+\Delta_\beta}\Bigr).
\end{multline*}

\emph{Step 3: the entropic term.} The volume of replica configurations with these overlaps, weighted by $e^{-\beta\lambda\sum_a\norm{\theta^a}_1}$, is written with conjugate variables scaled to survive $\beta\to\infty$: $\beta\hat Q/2$ for the norms, $\beta^2\hat\chi/2$ for the off-diagonal overlaps, $\beta\hat m$ for $m$. The identity $e^{\beta^2\hat\chi(\sum_ax_a)^2/2}=\E_ze^{\beta\sqrt{\hat\chi}z\sum_ax_a}$ decouples replicas and coordinates, and each coordinate sees the single-site weight $e^{-\beta[\hat Qx^2/2-x(\hat m\Theta_0+\sqrt{\hat\chi}Z)+\lambda\abs x]}$.

\emph{Step 4: zero temperature.} Let $\beta\to\infty$ with $\chi=\beta(Q-q)$ fixed and $k\to0$, keep the terms of order $\beta k$, and use Laplace's method (\cref{thm:toolkit:laplace}). With $\cE=Q-2m+\E\Theta_0^2$, the mean squared error,
\begin{equation}\label{eq:lasso:rs}
\begin{split}
e_0=\operatorname*{extr}\Bigl\{&-\frac{\hat QQ}2+\frac{\hat\chi\chi}2+\hat mm+\frac{\alpha\sigma^2+\cE}{2(1+\chi/\alpha)}\\
&+\E\min_{x\in\R}\Bigl[\frac{\hat Q}2x^2-x\bigl(\hat m\Theta_0+\sqrt{\hat\chi}Z\bigr)+\lambda\abs x\Bigr]\Bigr\},
\end{split}
\end{equation}
the extremum being over $Q,\chi,m,\hat Q,\hat\chi,\hat m$: the fraction is the residual energy, the minimum the entropic term of one coordinate.
\end{heuristic}

\begin{computation}{The saddle point of the replica potential}\label{comp:lasso:saddle}
\step \emph{Fields.} The potential is affine in $Q$ and $m$; stationarity gives $\hat Q=\hat m=\alpha/(\alpha+\chi)$. Stationarity in $\chi$ gives $\hat\chi=\alpha(\alpha\sigma^2+\cE)/(\alpha+\chi)^2$.
\step \emph{Single site.} As $\hat m=\hat Q$ and $\eta(cu;c\zeta)=c\,\eta(u;\zeta)$ for $c>0$, the minimizer is $\hat x=\eta(\Theta_0+\tau Z;\zeta)$ with $\tau=\sqrt{\hat\chi}/\hat Q$ and $\zeta=\lambda/\hat Q$.
\step \emph{Order parameters.} The derivative of a minimum is that at the minimizer, so stationarity in $\hat Q$, $\hat m$ gives $Q=\E\hat x^2$, $m=\E\hat x\Theta_0$, hence $\cE=\E(\hat x-\Theta_0)^2$, and step~1 gives $\tau^2=\hat\chi/\hat Q^2=\sigma^2+\cE/\alpha$: the first equation of \eqref{eq:lasso:fp}.
\step \emph{Susceptibility.} Stationarity in $\hat\chi$ gives $\chi=\E[\hat xZ]/\sqrt{\hat\chi}=s/\hat Q$, with $s=\Prob(\abs{\Theta_0+\tau Z}>\zeta)$, by Gaussian integration by parts (\cref{thm:toolkit:gaussianibp}); with step~1, $\hat Q=1-s/\alpha$.
\step \emph{Calibration.} So $\lambda=\zeta\hat Q=\zeta(1-s/\alpha)$, the second equation, and $\mathrm{MSE}=\cE=\alpha(\tau^2-\sigma^2)$. The same computation gives the training error $(\sigma^2+\cE/\alpha)/(1+\chi/\alpha)^2=(\tau\hat Q)^2$.
\end{computation}

This is prediction (i). Steps 1 and 4 of the computation also give $\chi=s/\hat Q=\alpha s/(\alpha-s)$: the susceptibility diverges as the fraction $s$ of selected coordinates approaches $\alpha$, that is, as the number of nonzero coordinates of $\hat\theta$ approaches the number of measurements. The prediction is reached through five uncontrolled steps: the continuation to $k\to0$ (\cref{ch:freeentropy}); the order of the limits $d\to\infty$, $k\to0$, $\beta\to\infty$; the \term{replica-symmetric ansatz}\index{replica symmetry!zero temperature}, supported by the convexity of $F$ and checked by the local stability of the solution, the \term{de Almeida--Thouless condition}\index{de Almeida--Thouless condition}, which \textcite{KabashimaEtAl2009} verified for noiseless $\ell_1$ recovery; the choice of extremum hidden in ``extr'', since the conjugate integrals run along imaginary axes; and the concentration of the overlaps.

\paragraph{The noiseless limit.} Take $\sigma=0$, $\lambda\to0$ and $\zeta=t\tau$. Exact recovery needs the map $\tau^2\mapsto\alpha^{-1}\E(\eta(\Theta_0+\tau Z;t\tau)-\Theta_0)^2$ to contract near $0$. As $\tau\to0$, $(\eta(\theta+\tau Z;t\tau)-\theta)/\tau\to Z-t\sgn\theta$ for $\theta\ne0$, and $\eta(\tau Z;t\tau)=\tau\eta(Z;t)$; so the slope of the map at $0$ is $J_\varepsilon(t)/\alpha$ for every law with $\Prob(\Theta_0\ne0)=\varepsilon$, and zero error is stable for some $t$ exactly when $\min_tJ_\varepsilon<\alpha$: prediction (ii), as \textcite{DonohoMalekiMontanari2009} derived it. For very sparse signals $\psi(\varepsilon)\sim2\varepsilon\log(1/\varepsilon)$ (\cref{ex:lasso:sparse}): about $2k\log(d/k)$ measurements recover $k$ nonzero entries. \Textcite{KabashimaEtAl2009} found the curve as the stability boundary of their ``success solution''.\sidenote{After submission they noticed that it coincides with the weak threshold of Donoho and Tanner's polytope computations, and added a note.}

\paragraph{Robust regression.} For $F(\theta)=\sum_i\rho(y_i-\ip{a_i}\theta)$, step~2 yields the Moreau envelope $\min_u\{(u-s)^2/(2c)+\rho(u)\}$, $c=\chi/\alpha$,\sidenote{By Laplace's method, $-\beta^{-1}\log$ of the convolution of $e^{-\beta\rho}$ with a Gaussian of variance $c/\beta$ tends to the Moreau envelope as $\beta\to\infty$.} and the saddle point is \eqref{eq:lasso:mest} with $r^2=\cE/\alpha$ (\cref{ex:lasso:robustreplica}): the susceptibility is the proximal step size. For $\rho(u)=u^2/2$ it gives the ridgeless risk $\sigma^2/(\alpha-1)$ of \cref{thm:ridge:risk}. For Huber's loss, $\frac12u^2$ for $\abs u\le k$ and $k\abs u-\frac12k^2$ otherwise, $\operatorname{prox}_c(x)$ is $x/(1+c)$ for $\abs x\le k(1+c)$ and $x-ck\sgn x$ otherwise.

\section{The stone}\label{sec:lasso:stone}

The rigorous route closest to the replica computation is Gordon's comparison in its convex form (\cref{ch:gordon}); the proof of \cref{thm:lasso:cgmt} follows the rows.

\begin{stone}{the replica computation}{Gordon's comparison (the CGMT)}
\stonestep{Replicate, or dualize}{Write $\min F$ as the ground state of $e^{-\beta F}$; compute $\E Z_\beta^k$ for $k$ replicas; $\beta\to\infty$ at the end.}{Write $\frac12\norm r^2=\max_u\{u^\top r-\frac12\norm u^2\}$: the LASSO is a min--max of a bilinear form in a Gaussian matrix. No replicas, no temperature: nothing corresponds to $k\to0$.}
\stonestep{Average over the design}{The fields $\ip{a_i}{\theta^a-\theta_0}$ are Gaussian with covariance given by the overlaps.}{Gordon's comparison replaces $u^\top G\tilde v$ by $\norm{\tilde v}g^\top u+\norm uh^\top\tilde v$ with independent Gaussian vectors $g$, $h$ (\cref{thm:gordon:cgmt}).}
\stonestep{The energetic term}{A Gaussian integral over the fields gives the residual energy $(\alpha\sigma^2+\cE)/(2(1+\chi/\alpha))$.}{Maximizing over the direction of $u$ leaves $b=\norm u$, the residual's norm over $\sqrt n$, and the term $b\sqrt{\norm v^2/n+\sigma^2}$.}
\stonestep{The entropic term}{A Gaussian field $\sqrt{\hat\chi}z$ decouples the coordinates; each solves $\min_x\frac{\hat Q}2x^2-x(\hat m\theta_{0j}+\sqrt{\hat\chi}z)+\lambda\abs x$.}{The identity $\sqrt x=\min_{\tau>0}\{\frac x{2\tau}+\frac\tau2\}$ separates the coordinates; each solves $\min_x\frac b{2\tau}(x-\theta_{0j})^2-bh_j(x-\theta_{0j})+\lambda\abs x$: the same, with $\hat Q=\hat m=b/\tau$, $\sqrt{\hat\chi}=b$, $z=h_j$.}
\stonestep{The saddle point}{Extremize over six order parameters (\cref{comp:lasso:saddle}); the method does not say which extremum.}{Eliminating $Q$ and $m$ turns the potential into $\alpha\Psi(b,\tau)$ (\cref{prop:lasso:rsao}); the auxiliary problem tends to $\max_b\min_\tau\Psi$, attained at $(b_*,\tau_*)$ only (\cref{lem:lasso:scalar}).}
\stonestep{Replica symmetry}{Assume one dominant overlap, supported by convexity and checked by stability.}{No ansatz. Convexity gives the matching upper bound (part (ii) of the CGMT) and, as strong convexity, a gap at every $v$ of the wrong norm: ``the overlap concentrates''.}
\stonestep{The decoupling principle}{The cavity field of coordinate $j$ is $\theta_{0j}$ plus Gaussian noise of variance $\tau^2$, for every statistic.}{AMP produces such fields with an exact state evolution (\cref{thm:amp:se}), and its iterates converge to the LASSO (\cref{thm:lasso:risk}).}
\stonestep{The noiseless limit}{The success solution $\tau=0$ is stable exactly when $\alpha>\psi(\varepsilon)$.}{The $\ell_1$ descent cone has statistical dimension $d\psi(s/d)$ up to $2\sqrt{d/s}$ (\cref{thm:lasso:statdim}); the kinematic formula gives the threshold (\cref{cor:lasso:dt}).}
\end{stone}

The proof does not merely confirm the replica answer; it reproduces the replica potential, and gives its conjugate fields meanings: $\sqrt{\hat\chi}$ is the residual's norm over $\sqrt n$, $1/\hat Q$ a proximal step size, the Hubbard--Stratonovich field Gordon's vector $h$. The continuation in $k$ and the limit $\beta\to\infty$ disappear because the comparison is made at zero temperature, and convexity decides the extremum. The routes part where convexity or Gaussianity fails: for $\ell_0$ the replica computation gives an unstable answer and Gordon's comparison one inequality, and for non-Gaussian designs the comparison has no counterpart.

\section{The proof: the LASSO risk by Gordon's comparison}\label{sec:lasso:proof}

The first proof analyzed an algorithm. Approximate message passing iterates $x^{t+1}=\eta(x^t+A^\top z^t;\zeta_t)$, $z^t=y-Ax^t+\alpha^{-1}z^{t-1}\langle\eta'\rangle$ (\cref{ch:amp}); by state evolution (\cref{thm:amp:se}; \textcite{BayatiMontanari2011}) the field $x^t+A^\top z^t$ behaves like $\theta_0+\tau_tZ$, with $\tau_{t+1}^2$ the right side of the first equation of \eqref{eq:lasso:fp} at $\tau_t$, and Bayati and Montanari proved that the iterates converge to the LASSO.

\begin{theorem}[The LASSO risk; \textcite{BayatiMontanari2012}, Theorem~1.5 and Corollary~1.6]\label{thm:lasso:risk}
Let $n/d\to\alpha\in(0,\infty)$, let $A$ have independent $\Normal(0,1/n)$ entries, and let $\theta_0\in\R^d$ and $w\in\R^n$ be deterministic, with empirical distributions converging weakly, with their second moments, to laws $P_{\Theta_0}$ and $P_W$ such that $\Prob(\Theta_0\ne0)>0$ and $\sigma^2=\E W^2>0$. Let $\lambda>0$. Then \eqref{eq:lasso:fp} has a unique solution $(\tau_*,\zeta_*)$ with $\tau_*,\zeta_*>0$, and for every $\varphi\colon\R^2\to\R$ with $\abs{\varphi(a)-\varphi(b)}\le L(1+\norm a+\norm b)\norm{a-b}$, almost surely,
\[
\lim_{d\to\infty}\frac1d\sum_{j=1}^d\varphi(\hat\theta_j,\theta_{0j})=\E\,\varphi\bigl(\eta(\Theta_0+\tau_*Z;\zeta_*),\Theta_0\bigr);
\]
in particular $\norm{\hat\theta-\theta_0}^2/d\to\alpha(\tau_*^2-\sigma^2)$ almost surely.
\end{theorem}

The noise need not be Gaussian; Bayati and Montanari remark that the Gaussian design is necessary for their technique. We prove the statement about the error by the second route, under two simplifications that keep the idea.\sidenote{For designs with a continuous distribution the LASSO minimizer is unique almost surely \parencite{Tibshirani2013}; the proof below does not need this, because it controls all minimizers at once.}

\begin{assumption}[Gaussian design and noise, independent signal]\label[assumption]{ass:lasso:model}
As $d\to\infty$, $n=n(d)$ with $n/d\to\alpha\in(0,\infty)$; $A$ has independent $\Normal(0,1/n)$ entries; $w$ has independent $\Normal(0,\sigma^2)$ entries, $\sigma>0$; $\theta_0$ has independent entries with a law $P_{\Theta_0}$, $\E\Theta_0^2<\infty$; $A$, $w$, $\theta_0$ are independent; $\lambda>0$ is fixed.
\end{assumption}

\subsection{The scalar problem}\label{sec:lasso:scalar}

The limit of Gordon's auxiliary problem is a convex problem over random variables, whose minimizer is the decoupled estimate and whose first-order conditions are \eqref{eq:lasso:fp}.

\begin{lemma}[The scalar problem]\label[lemma]{lem:lasso:scalar}
Let $\sigma,\lambda,\alpha>0$, $\E\Theta_0^2<\infty$, $Z\sim\Normal(0,1)$ independent of $\Theta_0$, and let $\cH$ be the Hilbert space of square-integrable functions of $(\Theta_0,Z)$, with $\norm V^2=\E V^2$. For $V\in\cH$ and $b,\tau>0$ put
\begin{gather}
M(V)=\sqrt{\sigma^2+\norm V^2/\alpha}-\E[VZ]/\alpha,\notag\\
\cF(V)=\tfrac12M(V)_+^2+\tfrac\lambda\alpha\E\abs{\Theta_0+V},\label{eq:lasso:F}\\
\Psi(b,\tau)=\frac{b\sigma^2}{2\tau}+\frac{b\tau}2-\frac{b^2}2+\frac1\alpha\E q_b(\tau;\Theta_0,Z),\label{eq:lasso:Psi}\\
q_b(\tau;\theta,z)=\min_{x\in\R}\Bigl\{\frac b{2\tau}(x-\theta)^2-bz(x-\theta)+\lambda\abs x\Bigr\}.\notag
\end{gather}
\begin{enumerate}
\item $\cF$ has a unique minimizer $V_*$, and $b_*\defeq M(V_*)>0$.
\item With $\tau_*=\sqrt{\sigma^2+\norm{V_*}^2/\alpha}$ and $\zeta_*=\lambda\tau_*/b_*$: $\Theta_0+V_*=\eta(\Theta_0+\tau_*Z;\zeta_*)$ almost surely, $(\tau_*,\zeta_*)$ is the only solution of \eqref{eq:lasso:fp} with $\tau,\zeta>0$, and $b_*=\tau_*(1-\Prob(\abs{\Theta_0+\tau_*Z}>\zeta_*)/\alpha)$.
\item $\Psi(b_*,\cdot)$ is strictly convex on $(0,\infty)$, with unique minimizer $\tau_*$ and minimum $J_*\defeq\cF(V_*)=\frac12b_*^2+\frac\lambda\alpha\E\abs{\eta(\Theta_0+\tau_*Z;\zeta_*)}$, and $b_*$ is the only maximizer of $b\mapsto\inf_{\tau>0}\Psi(b,\tau)$ on $(0,\infty)$.
\end{enumerate}
\end{lemma}

\begin{proof}
The plan is to obtain $V_*$ from convexity and coercivity, to read the fixed point off its first-order conditions, and to recover $\Psi$ by undoing $\sqrt x=\min_\tau\{\frac x{2\tau}+\frac\tau2\}$.

\emph{Part 1.} The map $V\mapsto\sqrt{\sigma^2+\norm V^2/\alpha}$ is the norm of the affine map $V\mapsto(V/\sqrt\alpha,\sigma)$, so $M$, and with it $\cF$, is convex and continuous. Choose $b_0>0$ with $D(\lambda/b_0)\le\alpha/4$. The bounds $\frac12m_+^2\ge b_0m-\frac12b_0^2$ and $VZ-\frac\lambda{b_0}\abs V\le\abs V(\abs Z-\frac\lambda{b_0})_+$, and the Cauchy--Schwarz inequality, give $\cF(V)\ge\frac{b_0}{2\sqrt\alpha}\norm V-\frac12b_0^2-\frac\lambda\alpha\E\abs{\Theta_0}$. A convex, continuous, coercive function on a Hilbert space attains its minimum, since its sublevel sets are weakly compact and it is weakly lower semicontinuous. If a minimizer $V$ had $M(V)\le0$, then $V\ne-\Theta_0$, because $M(-\Theta_0)=\sqrt{\sigma^2+\E\Theta_0^2/\alpha}>0$, so $\cF(V)=\frac\lambda\alpha\E\abs{\Theta_0+V}>0$; the points $V_s=(1-s)V-s\Theta_0$ have $\E\abs{\Theta_0+V_s}=(1-s)\E\abs{\Theta_0+V}$ and $M(V_s)\le sM(-\Theta_0)$, so $\cF(V_s)\le(1-s)\cF(V)+\frac12s^2M(-\Theta_0)^2<\cF(V)$ for small $s>0$, a contradiction. Two distinct minimizers would give, at their midpoint, a value of $M$ below the average of theirs, because equality in the triangle inequality forces positively proportional vectors and $(V/\sqrt\alpha,\sigma)$ has the fixed coordinate $\sigma>0$; as $m\mapsto\frac12m_+^2$ is convex and increasing on $[0,\infty)$, $\cF$ would be smaller there than its minimum.

\emph{Part 2.} Let $X_*=\Theta_0+V_*$ and $T=b_*(V_*/\tau_*-Z)$. Near $V_*$ the function $\frac12M_+^2=\frac12M^2$ has gradient $T/\alpha$, and the one-sided derivative of $s\mapsto\E\abs{X_*+s\Delta}$ at $0$ is $\E[\sgn(X_*)\Delta;X_*\ne0]+\E[\abs\Delta;X_*=0]$. Minimality makes the one-sided derivative of $\cF$ nonnegative in every direction $\Delta$; the directions $-(T+\lambda\sgn X_*)\ind\{X_*\ne0\}$ and $\pm\ind_B$ with $B\subset\{X_*=0\}$ give $T=-\lambda\sgn X_*$ on $\{X_*\ne0\}$ and $\abs T\le\lambda$ on $\{X_*=0\}$. Since $T=(b_*/\tau_*)(X_*-\Theta_0-\tau_*Z)$, these are the conditions for $X_*=\eta(\Theta_0+\tau_*Z;\zeta_*)$. The definition of $\tau_*$ is the first equation of \eqref{eq:lasso:fp}. Conditionally on $\Theta_0$, $z\mapsto\eta(\Theta_0+\tau_*z;\zeta_*)$ is Lipschitz with derivative $\tau_*\ind\{\abs{\Theta_0+\tau_*z}>\zeta_*\}$ almost everywhere, so Gaussian integration by parts (\cref{thm:toolkit:gaussianibp}) gives $\E[V_*Z]=\tau_*\Prob(\abs{\Theta_0+\tau_*Z}>\zeta_*)$; this is the formula for $b_*=M(V_*)$, and $\lambda=\zeta_*b_*/\tau_*$ is the second equation. Conversely, if $(\tau,\zeta)$ with $\tau,\zeta>0$ solves \eqref{eq:lasso:fp}, then $V=\eta(\Theta_0+\tau Z;\zeta)-\Theta_0$ has $M(V)=b\defeq\lambda\tau/\zeta$ by the same integration by parts, and $T=(b/\tau)(V-\tau Z)=-\lambda S$ for a pointwise subgradient $S$ of $\abs\cdot$ at $\Theta_0+V$. The gradient inequality for $\frac12M_+^2$ and the subgradient inequality for $\abs\cdot$ give $\cF(V')\ge\cF(V)+\frac1\alpha\E[(T+\lambda S)(V'-V)]=\cF(V)$ for all $V'$, so $V=V_*$ and $(\tau,\zeta)=(\tau_*,\zeta_*)$.

\emph{Part 3.} For fixed $(\theta,z)$ the function $(x,\tau)\mapsto\frac b{2\tau}(x-\theta)^2$ is jointly convex (a perspective), so $q_b(\cdot;\theta,z)$ and its expectation are convex in $\tau$, and $b\sigma^2/(2\tau)$ is strictly convex: $\Psi(b_*,\cdot)$ is strictly convex. The minimizer in $q_b$ is $\eta(\theta+\tau z;\lambda\tau/b)$, square integrable at $(\Theta_0,Z)$, so with $V=x-\Theta_0$
\[
\Psi(b_*,\tau)=\min_{V\in\cH}\Bigl\{b_*\Bigl(\frac{\sigma^2+\norm V^2/\alpha}{2\tau}+\frac\tau2\Bigr)-\frac{b_*}\alpha\E[VZ]+\frac\lambda\alpha\E\abs{\Theta_0+V}\Bigr\}-\frac{b_*^2}2.
\]
The arithmetic--geometric mean inequality bounds the inner parenthesis below by $\sqrt{\sigma^2+\norm V^2/\alpha}$, so $\Psi(b_*,\tau)\ge\min_VG(V)$ with $G(V)=b_*M(V)-\frac12b_*^2+\frac\lambda\alpha\E\abs{\Theta_0+V}$. The convex function $b_*M$ has gradient $T/\alpha$ at $V_*$, as $\frac12M^2$ has, so the subgradient argument of part~2, with $S=-T/\lambda$, shows that $V_*$ minimizes $G$; and $G(V_*)=\frac12b_*^2+\frac\lambda\alpha\E\abs{X_*}=\cF(V_*)$. At $\tau=\tau_*$ the inequality is an equality at $V=V_*$, which is the pointwise minimizer there because $\lambda\tau_*/b_*=\zeta_*$; so $\Psi(b_*,\tau_*)=\cF(V_*)$, and $\tau_*$ is the unique minimizer. For $b>0$ the display holds with $b$ for $b_*$, and minimizing over $\tau$ first gives $\inf_\tau\Psi(b,\tau)=\inf_VG_b(V)\le G_b(V_*)=bb_*-\frac12b^2+\frac\lambda\alpha\E\abs{X_*}$, where $G_b$ is $G$ with $b$ for $b_*$; this is below $J_*$ unless $b=b_*$.
\end{proof}

\subsection{Gordon's comparison for the LASSO}\label{sec:lasso:cgmt}

The Gaussian noise can be absorbed into the matrix.\sidenote{The device is used by \textcite{ThrampoulidisOymakHassibi2015} and \textcite{MiolaneMontanari2021}. Without it the auxiliary problem contains $\norm{\norm vg/\sqrt n-w}$, which is not convex in $v$.} Let $\tilde G=[\sqrt nA,\,w/\sigma]\in\R^{n\times(d+1)}$, which has independent $\Normal(0,1)$ entries, independent of $\theta_0$, and $\tilde v=(-v,\sigma\sqrt n)$ for $v\in\R^d$. Then $\tilde G\tilde v=\sqrt n(w-Av)$, and
\begin{equation}\label{eq:lasso:dual}
\max_{u\in\R^n}\Bigl\{u^\top\tilde G\tilde v-\frac n2\norm u^2\Bigr\}=\frac12\norm{w-Av}^2=F(\theta_0+v)-\lambda\norm{\theta_0+v}_1,
\end{equation}
attained at $u=(w-Av)/\sqrt n$. For compact $S\subset\R^d$, $C>0$ and $\Upsilon(v,u)=-\frac n2\norm u^2+\lambda\norm{\theta_0+v}_1$ define the primary problem and the \term{auxiliary optimization}\index{auxiliary optimization}
\begin{align*}
\cP(S)&=\min_{v\in S}\max_{\norm u\le C}\bigl\{u^\top\tilde G\tilde v+\Upsilon(v,u)\bigr\},\\
\cA(S)&=\min_{v\in S}\max_{\norm u\le C}\bigl\{\norm{\tilde v}g^\top u+\norm u\,\tilde h^\top\tilde v+\Upsilon(v,u)\bigr\},
\end{align*}
with $g\in\R^n$ and $\tilde h=(h,h_0)\in\R^d\times\R$ independent standard Gaussian vectors, independent of $\theta_0$. The convex Gaussian min-max theorem (\cref{thm:gordon:cgmt}), in the form of Theorem~6.1 of \textcite{ThrampoulidisAbbasiHassibi2018}, whose parts (i) and (ii) are those of Theorem~3 of \textcite{ThrampoulidisOymakHassibi2015}, applies conditionally on $\theta_0$ to the compact set $\{\tilde v:v\in S\}$, and gives, for every $c\in\R$:
\textup{(i)} $\Prob(\cP(S)<c)\le2\Prob(\cA(S)\le c)$, which is Gordon's inequality \parencite{Gordon1985,Gordon1988} and needs no convexity; and
\textup{(ii)} if $S$ is convex, $\Prob(\cP(S)>c)\le2\Prob(\cA(S)\ge c)$, because $\Upsilon$ is convex in $v$ and concave in $u$.
In $\cA(S)$ the maximum over the direction of $u$ at fixed $\norm u=b$ puts $u$ along $g$, and $\norm{\tilde v}=\sqrt n\sqrt{\norm v^2/n+\sigma^2}$, so
\begin{equation}\label{eq:lasso:lb}
\frac{\cA(S)}n=\min_{v\in S}\max_{0\le b\le C}\ell_b(v),\qquad\ell_b(v)=bX(v)-\frac{b^2}2+\frac\lambda n\norm{\theta_0+v}_1,
\end{equation}
with $X(v)=\sqrt{\norm v^2/n+\sigma^2}\,\norm g/\sqrt n-\ip hv/n+\sigma h_0/\sqrt n$. As every claim below concerns laws, we take the entries of $\theta_0$, $g$, $h$ from fixed infinite independent sequences, so that the strong law of large numbers applies.

\begin{lemma}[A priori bounds]\label[lemma]{lem:lasso:bounded}
Assume \cref{ass:lasso:model}. There is $K_0$, depending on $\alpha,\sigma,\lambda,\E\abs{\Theta_0}$, such that for $K\ge K_0$ and $C\ge2\sigma+(3+\alpha^{-1/2})K$, with probability tending to one, every minimizer of $F$ lies in $\{\theta:\norm{\theta-\theta_0}<K\sqrt n\}$, and $\norm{w-Av}\le C\sqrt n$ whenever $\norm v\le K\sqrt n$.
\end{lemma}

\begin{proof}
The plan is to bound $A$ and $w$, and to show by Gordon's inequality that $F$ is large on the sphere of radius $K\sqrt n$ around $\theta_0$; convexity then traps the minimizers inside. By the Sudakov--Fernique inequality (\cref{thm:toolkit:slepian}; \textcite[\S7.3]{Vershynin2018}) $\E\norm{\sqrt nA}_{\mathrm{op}}\le\sqrt n+\sqrt d$, and Gaussian concentration (\cref{thm:toolkit:lipschitz}) gives $\norm A_{\mathrm{op}}\le2+\sqrt{d/n}\le3+\alpha^{-1/2}$ with probability tending to one; with $\norm w\le2\sigma\sqrt n$ this gives the second claim.

Let $F_0=\frac12\sigma^2+\frac\lambda\alpha\E\abs{\Theta_0}$, the limit of $F(\theta_0)/n$; take $b_0\in(0,1]$ with $D(\lambda/b_0)\le\alpha/4$ and $K_0=\frac2{b_0}(F_0+3+\frac12b_0^2+\frac\lambda\alpha\E\abs{\Theta_0})$, so that $b_0\le1\le C$. On $S_K=\{v:\norm v=K\sqrt n\}$, the bound $\lambda\abs{v_j}-b_0h_jv_j\ge-b_0(\abs{h_j}-\lambda/b_0)_+\abs{v_j}$ and the Cauchy--Schwarz inequality give
\[
\ell_{b_0}(v)\ge b_0K\frac{\norm g}{\sqrt n}-\frac{b_0\sigma\abs{h_0}}{\sqrt n}-\frac{b_0^2}2-\frac\lambda n\norm{\theta_0}_1-b_0K\frac{\norm{(\abs h-\lambda/b_0)_+}}{\sqrt n}.
\]
By the law of large numbers the right side tends to a limit at least $\frac12b_0K-\frac12b_0^2-\frac\lambda\alpha\E\abs{\Theta_0}\ge F_0+3$, since $\norm{(\abs h-\lambda/b_0)_+}^2/n\to D(\lambda/b_0)/\alpha\le\frac14$. So $\cA(S_K)/n\ge F_0+2$ with probability tending to one, and part (i) of the CGMT, which needs only compactness, gives $\Prob(\cP(S_K)/n<F_0+2)\to0$. Restricting $u$ lowers the maximum in \eqref{eq:lasso:dual}, so $F(\theta_0+v)\ge\cP(S_K)$ on $S_K$, and with probability tending to one $\min_{S_K}F(\theta_0+\cdot)>F(\theta_0)$. Then no minimizer $\hat\theta$ has $\norm{\hat\theta-\theta_0}\ge K\sqrt n$: the segment from $\theta_0$ to $\hat\theta$ would meet the sphere at a point where, by convexity, $F\le\max\{F(\theta_0),F(\hat\theta)\}=F(\theta_0)$.
\end{proof}

\begin{theorem}[The LASSO risk by Gordon's comparison]\label{thm:lasso:cgmt}
Assume \cref{ass:lasso:model}, and let $\tau_*,\zeta_*,b_*,J_*$ be as in \cref{lem:lasso:scalar}. For every $\gamma>0$ the probability that some minimizer $\hat\theta$ of $F$ has $\abs{\norm{\hat\theta-\theta_0}^2/d-\alpha(\tau_*^2-\sigma^2)}\ge\gamma$ tends to zero, and $\min F/n\to J_*=\frac12b_*^2+\frac\lambda\alpha\E\abs{\eta(\Theta_0+\tau_*Z;\zeta_*)}$ in probability.
\end{theorem}

\begin{proof}
The plan is to show that the auxiliary problem tends to $J_*$, with an upper bound from a candidate built on the fixed point and a lower bound from fixing $b=b_*$, and that it exceeds $J_*$ by a fixed amount where the norm of $v$ is wrong; the CGMT transfers both facts.

\emph{Step 1: reduction.} Let $a_*=\sqrt{\tau_*^2-\sigma^2}$; as $n/d\to\alpha$, it suffices to control $\norm{\hat\theta-\theta_0}/\sqrt n-a_*$. Fix $\gamma>0$, $K\ge\max\{K_0,a_*+1\}$, $C\ge\max\{b_*+1,2\sigma+(3+\alpha^{-1/2})K\}$, the ball $B=\{v:\norm v\le K\sqrt n\}$ and the compact set $S^c=\{v\in B:\abs{\norm v/\sqrt n-a_*}\ge\gamma\}$. On the event $\Omega_1$ of \cref{lem:lasso:bounded} the constraint $\norm u\le C$ does not bind in \eqref{eq:lasso:dual} on $B$, so $\cP(B)=\min F$, $\cP(S^c)=\min_{S^c}F(\theta_0+\cdot)$, and a minimizer with $\hat\theta-\theta_0\in S^c$ forces $\cP(S^c)\le\cP(B)$. It suffices to show $\Prob(\cP(S^c)\le\cP(B))\to0$.

\emph{Step 2: upper bound.} Let $x^*_j=\eta(\theta_{0j}+\tau_*h_j;\zeta_*)$, $v^*=x^*-\theta_0$ and $X_*=\Theta_0+V_*$. By the strong law of large numbers and integration by parts, almost surely, $\norm{v^*}^2/n\to a_*^2<K^2$, $\ip h{v^*}/n\to\tau_*\Prob(\abs{\Theta_0+\tau_*Z}>\zeta_*)/\alpha$ and $\frac\lambda n\norm{x^*}_1\to\frac\lambda\alpha\E\abs{X_*}$; so $X(v^*)\to b_*$ by \cref{lem:lasso:scalar}, and
\[
\limsup_d\cA(B)/n\le\limsup_d\max_{b\ge0}\ell_b(v^*)=\tfrac12b_*^2+\tfrac\lambda\alpha\E\abs{X_*}=J_*\qquad\text{almost surely.}
\]

\emph{Step 3: lower bound.} As $b_*\le C$, $\cA(S)/n\ge\min_S\ell_{b_*}$ for $S\subset B$. With $\omega_d=\norm g/\sqrt n$, the identity $\omega_d\sqrt x=\min_{\tau>0}\{\frac x{2\tau}+\frac{\omega_d^2\tau}2\}$ and completing squares coordinatewise give $\min_{\R^d}\ell_{b_*}=\min_{\tau>0}\Gamma_d(\tau)$, with
\[
\Gamma_d(\tau)=\frac{b_*\sigma^2}{2\tau}+\frac{b_*\omega_d^2\tau}2+\frac{b_*\sigma h_0}{\sqrt n}-\frac{b_*^2}2+\frac1n\sum_{j=1}^dq_{b_*}(\tau;\theta_{0j},h_j),
\]
with $q_b$ as in \cref{lem:lasso:scalar}; at fixed $\tau$ the minimum over $v$ is attained at $v_{\tau,j}=\eta(\theta_{0j}+\tau h_j;\lambda\tau/b_*)-\theta_{0j}$. Each $q_{b_*}(\cdot;\theta,z)$ is convex (part~3 of the proof of \cref{lem:lasso:scalar}), and $\abs{q_{b_*}}\le\lambda\abs\theta+\frac12b_*\tau z^2$, so $\Gamma_d$ is convex and $\Gamma_d(\tau)\to\Psi(b_*,\tau)$ almost surely, simultaneously for all rational $\tau>0$; convex functions converging on a dense set converge uniformly on compacts \parencite[Theorem~10.8]{Rockafellar1970}. As $\Psi(b_*,\tau_*\pm\delta)>J_*$ for $\delta\in(0,\tau_*)$, for large $d$ $\Gamma_d(\tau_*\pm\delta)>\Gamma_d(\tau_*)$, so by convexity $\Gamma_d$ attains its minimum, at points $\hat\tau_d\in(\tau_*-\delta,\tau_*+\delta)$. Hence, almost surely, $\hat\tau_d\to\tau_*$ and $\min\Gamma_d\to J_*$. The strictly convex $\ell_{b_*}$ has the unique minimizer $\hat v=v_{\hat\tau_d}$, and since $\eta$ is $1$-Lipschitz in both arguments, $\norm{v_\tau-v_{\tau'}}\le\abs{\tau-\tau'}(\norm h+\lambda\sqrt d/b_*)$; with $v_{\tau_*}=v^*$ this gives $\norm{\hat v}/\sqrt n\to a_*$ almost surely.

\emph{Step 4: the gap.} On $B$, $f(v)=\sqrt{\norm v^2/n+\sigma^2}$ has Hessian $\frac1{nf}(I-\frac{vv^\top}{nf^2})\succeq\frac{\sigma^2}{nf^3}I\succeq\frac{\sigma^2}{n(K^2+\sigma^2)^{3/2}}I$; with $\mu=b_*\sigma^2/(2(K^2+\sigma^2)^{3/2})$ and $\omega_d\ge\frac12$, $\ell_{b_*}-\frac\mu{2n}\norm\cdot^2$ is convex on $B$. For $\hat v,v\in B$ and $s\in(0,1]$, convexity along the segment gives $\ell_{b_*}(\hat v)\le\ell_{b_*}(\hat v+s(v-\hat v))\le(1-s)\ell_{b_*}(\hat v)+s\ell_{b_*}(v)-\frac\mu{2n}s(1-s)\norm{v-\hat v}^2$, and $s\to0$ gives $\ell_{b_*}(v)\ge\ell_{b_*}(\hat v)+\frac\mu{2n}\norm{v-\hat v}^2$. For large $d$, $\hat v\in B$ and $\abs{\norm{\hat v}/\sqrt n-a_*}<\gamma/2$, so each $v\in S^c$ has $\norm{v-\hat v}\ge\frac\gamma2\sqrt n$, and with $\Delta=\mu\gamma^2/8$
\[
\cA(S^c)/n\ge\min_{S^c}\ell_{b_*}\ge\min\Gamma_d+\Delta\to J_*+\Delta.
\]

\emph{Step 5: transfer.} Part (i) of the CGMT and step~4 give $\Prob(\cP(S^c)/n<J_*+\frac\Delta2)\le2\Prob(\cA(S^c)/n\le J_*+\frac\Delta2)\to0$; part (ii), on the convex set $B$, and step~2 give $\Prob(\cP(B)/n>J_*+\frac\Delta4)\le2\Prob(\cA(B)/n\ge J_*+\frac\Delta4)\to0$. So $\Prob(\cP(S^c)\le\cP(B))\to0$, and step~1 proves the first claim. For the second, steps 2 and 3 give $\cA(B)/n\to J_*$ in probability, parts (i) and (ii) give $\Prob(\abs{\cP(B)/n-J_*}>t)\le2\Prob(\abs{\cA(B)/n-J_*}\ge t)\to0$, and $\cP(B)=\min F$ on $\Omega_1$.
\end{proof}

The same gap controls every function of $v$ that is Lipschitz for $\norm\cdot/\sqrt n$ and converges at $\hat v$, such as the overlap $m$ (\cref{ex:lasso:overlap}); the empirical distribution of the $(\hat\theta_j,\theta_{0j})$ needs a stability argument \parencite{MiolaneMontanari2021}.

\subsection{The replica potential is the auxiliary problem}\label{sec:lasso:rsao}

\begin{proposition}[The replica potential and $\Psi$]\label[proposition]{prop:lasso:rsao}
Let $\Phi_{\mathrm{RS}}$ be the function in braces in \eqref{eq:lasso:rs}. Stationarity in $Q$ and $m$ holds exactly when $\hat Q=\hat m=\alpha/(\alpha+\chi)$, and under these relations the change of variables $\hat Q=b/\tau$, $\hat\chi=b^2$, $\chi=\alpha(\tau/b-1)$, with $b,\tau>0$, gives $\Phi_{\mathrm{RS}}=\alpha\Psi(b,\tau)$.
\end{proposition}

\begin{proof}
The plan is to substitute. The terms containing $Q$ or $m$ are $-\frac12\hat QQ+\hat mm+(Q-2m)/(2(1+\chi/\alpha))$; their derivatives vanish exactly when $\hat Q=\hat m=1/(1+\chi/\alpha)$, and then they cancel, leaving $\frac12\hat\chi\chi+\frac12\hat Q(\alpha\sigma^2+\E\Theta_0^2)$ plus the expected minimum. With $\hat Q=b/\tau$ and $\sqrt{\hat\chi}=b$ the bracket is $\frac b{2\tau}(x-\Theta_0)^2-bZ(x-\Theta_0)+\lambda\abs x-\frac b{2\tau}\Theta_0^2-bZ\Theta_0$, with $\E[Z\Theta_0]=0$, and $\frac12\hat\chi\chi=\frac12\alpha b(\tau-b)$. The terms in $\E\Theta_0^2$ cancel, and what remains is $\frac12\alpha b\tau-\frac12\alpha b^2+\frac{\alpha b\sigma^2}{2\tau}+\E\min_x[\cdots]=\alpha\Psi(b,\tau)$.
\end{proof}

With the computation's $\hat Q=1-s/\alpha$, the relation $\hat Q=b/\tau$ is the formula for $b_*$, and the replica training error $(\tau\hat Q)^2$ is $b_*^2$, the squared norm of the dual variable, the residual over $\sqrt n$. The proof uses no property of the penalty, so the identification holds for every separable penalty, the ridge penalty of \cref{ch:ridge} included.

\subsection{What the general proofs add}\label{sec:lasso:general}

\Cref{thm:lasso:risk} gives every pseudo-Lipschitz statistic, for deterministic signals and general noise. \Textcite{ThrampoulidisAbbasiHassibi2018} treat general convex losses and separable penalties by our route (their Theorem~4.1 and \S5.4). \Textcite{MiolaneMontanari2021} prove Wasserstein convergence of the empirical distribution uniformly over $\lambda\in[\lambda_{\min},\lambda_{\max}]$ and over $\ell_p$ balls; their max--min (8) is our $\Psi$ up to a constant, and their Proposition~3.1 our \cref{lem:lasso:scalar}; general Gaussian designs are treated by \textcite{CelentanoMontanariWei2023}.

For robust regression with design $x_i\iid\Normal(0,I_p)$ and $n/p\to\alpha>1$, the system \eqref{eq:lasso:mest} gives the limit of $\norm{\hat\beta-\beta^*}^2$ by AMP for strongly convex smooth losses \parencite[Theorem~4.1]{DonohoMontanari2016}, by the CGMT for convex losses under conditions that Huber's loss meets even with Cauchy noise \parencite[Theorem~4.1, \S5.1]{ThrampoulidisAbbasiHassibi2018}, and, for designs with independent non-Gaussian entries and strongly convex losses, by leave-one-out arguments in the unrefereed preprint of \textcite{ElKaroui2013}, who notes that such ideas ``seem to be known in Physics under the name `cavity method'\,''; its refereed revision \parencite{ElKaroui2018} keeps only ridge-regularized estimators, with elliptical-like designs, as does the universality theorem of \textcite{HanShen2023}. Donoho and Montanari also show that the system forces $R\ge1/((\alpha-1)I(P_W))$, $I$ the Fisher information of the noise: a \term{variance inflation}\index{variance inflation} that no loss avoids (\cref{ex:lasso:fisher}).

\section{The proof: the Donoho--Tanner threshold as a statistical dimension}\label{sec:lasso:dt}

Without noise, the LASSO tends as $\lambda\to0$ to $\ell_1$ minimization, $\min\norm\theta_1$ subject to $A\theta=y$, which succeeds if $\theta_0$ is its unique solution. The \term{descent cone}\index{descent cone} of a convex $f$ at $x$ is $\cD(f,x)=\{v:f(x+sv)\le f(x)\text{ for some }s>0\}$, and $\ell_1$ minimization succeeds exactly when $\ker A\cap\cD(\norm\cdot_1,\theta_0)=\{0\}$: a nonzero $v$ in the intersection makes $\theta_0+sv$ a competitor, and a competitor $\theta$ puts $\theta-\theta_0$ in the intersection. The \term{statistical dimension}\index{statistical dimension} of a closed convex cone $\cC\subset\R^d$ is $\delta(\cC)=\E\norm{\Pi_\cC(g)}^2$, with $g\sim\Normal(0,I_d)$ and $\Pi_\cC$ the projection, and that of a convex cone is that of its closure; a $k$-dimensional subspace has $\delta=k$, and the orthant $\delta=d/2$, which gives the prologue's number two.\sidenote{Up to one, $\delta(\cC)$ is the squared Gaussian width of $\cC\cap S^{d-1}$ \parencite[Proposition~10.2]{AmelunxenEtAl2014}, the quantity of Gordon's escape-through-a-mesh argument.}

\begin{theorem*}[\textcite{AmelunxenEtAl2014}, Theorem~II]
Fix $p\in(0,1)$, $a_p=\sqrt{8\log(4/p)}$, a fixed $x_0\in\R^d$, a proper convex $f$, and $A\in\R^{n\times d}$ with independent standard normal entries. If $n\le\delta(\cD(f,x_0))-a_p\sqrt d$, then $\min f(x)$ subject to $Ax=Ax_0$ succeeds with probability at most $p$; if $n\ge\delta(\cD(f,x_0))+a_p\sqrt d$, with probability at least $1-p$.
\end{theorem*}

Its proof combines the optimality condition with an approximate \term{kinematic formula}\index{kinematic formula} for a random subspace meeting a cone; what is particular to $\ell_1$ is the statistical dimension.

\begin{theorem}[The $\ell_1$ descent cone; \textcite{AmelunxenEtAl2014}, Proposition~4.5]\label{thm:lasso:statdim}
If $x\in\R^d$ has $s$ nonzero entries, $1\le s\le d-1$, and $\cD=\cD(\norm\cdot_1,x)$, then $d\,\psi(s/d)-2\sqrt{d/s}\le\delta(\cD)\le d\,\psi(s/d)$.
\end{theorem}

\begin{proof}
The plan, which is the argument of \textcite[Proposition~4.1 and Appendix~C.2]{AmelunxenEtAl2014} specialized to $\ell_1$, is to write $\delta(\cD)$ as an expected squared distance to the cone generated by the subdifferential, to exchange expectation and minimum for the upper bound, and to bound the cost of the exchange by two variances. The subdifferential $S=\{u:u_i=\sgn x_i\text{ if }x_i\ne0,\ \abs{u_i}\le1\text{ otherwise}\}$, and with it $\cD$, depends only on the signs of $x$, and signed permutations preserve the law of $g$; so we take $x=(1,\dots,1,0,\dots,0)$ with $s$ ones, for which $\ip ux=s$ on $S$.

By Moreau's decomposition $\norm{\Pi_\cC(g)}=\operatorname{dist}(g,\cC^\circ)$, and the polar of the descent cone at a point where the subdifferential is compact and does not contain $0$ is the cone it generates \parencite[Corollary~23.7.1]{Rockafellar1970}, here the closed cone $\mathcal K=\bigcup_{t\ge0}tS$. Computing the nearest point of $tS$ coordinatewise,
\[
\delta(\cD)=\E\min_{t\ge0}F_g(t),\qquad F_g(t)=\operatorname{dist}^2(g,tS)=\sum_{i\le s}(g_i-t)^2+\sum_{i>s}(\abs{g_i}-t)_+^2.
\]
Since $\E F_g(t)=s(1+t^2)+(d-s)D(t)=dJ_{s/d}(t)$, exchanging expectation and minimum gives $\delta(\cD)\le d\psi(s/d)$.

For the lower bound write $J=J_{s/d}$. From $D(t)=2[(1+t^2)\Phi(-t)-t\phi(t)]$, $D'(t)=-4[\phi(t)-t\Phi(-t)]$ and $D''(t)=4\Phi(-t)>0$; so $J$ is strictly convex with $J'(0)<0$ and $J'(\infty)=\infty$, and has a unique minimizer $t_\star>0$. Each $F_g$ is convex and differentiable, $F_g'(t)=2st-2\sum_{i\le s}g_i-2\sum_{i>s}(\abs{g_i}-t)_+$, and $\E F_g'(t)=dJ'(t)$. With $t_g$ the minimizer of $F_g$ on $[0,\infty)$, the tangent at $t_\star$ gives $F_g(t_g)\ge F_g(t_\star)+(t_g-t_\star)F_g'(t_\star)$, and since $\E F_g'(t_\star)=0$,
\[
\delta(\cD)\ge d\psi(s/d)+\Cov\bigl(t_g,F_g'(t_\star)\bigr)\ge d\psi(s/d)-\sqrt{\Var(t_g)\Var(F_g'(t_\star))}.
\]
The Gaussian Poincar\'e inequality (\cref{thm:toolkit:poincare}) gives $\Var f(g)\le L^2$ for $L$-Lipschitz $f$. The gradient of $g\mapsto F_g'(t_\star)$ has entries $-2$ and $-2\sgn(g_i)\ind\{\abs{g_i}>t_\star\}$, so $\Var F_g'(t_\star)\le4d$. The nearest point of $\mathcal K$ to $g$ is that of $t_gS$, by uniqueness of projections, so $t_g=\ip{\Pi_{\mathcal K}(g)}x/s$; as $\Pi_{\mathcal K}$ is $1$-Lipschitz, $\abs{t_g-t_{g'}}\le\norm x\norm{g-g'}/s=\norm{g-g'}/\sqrt s$ and $\Var t_g\le1/s$. Hence $\delta(\cD)\ge d\psi(s/d)-2\sqrt{d/s}$.
\end{proof}

For a fixed signal this gives the \term{weak threshold}\index{weak threshold} of Donoho and Tanner, as opposed to the strong threshold, at which all $k$-sparse vectors are recovered at once.

\begin{corollary}[The weak threshold]\label[corollary]{cor:lasso:dt}
Let $s/d\to\varepsilon\in(0,1)$ and $n/d\to\alpha\in(0,1)$, let $\theta_0$ have $s$ nonzero entries, and let $A$ be Gaussian. If $\alpha>\psi(\varepsilon)$, $\ell_1$ minimization succeeds with probability tending to one; if $\alpha<\psi(\varepsilon)$, it fails with probability tending to one. The function $\psi$ is continuous and increasing from $\psi(0)=0$ to $\psi(1)=1$, so the threshold is $\varepsilon_c(\alpha)=\psi^{-1}(\alpha)$.
\end{corollary}

\begin{proof}
The plan is to insert \cref{thm:lasso:statdim} into Theorem~II. For each $t$, $J_\varepsilon(t)$ is affine in $\varepsilon$ with slope $1+t^2-D(t)$, positive for $t>0$; the minimizing $t$ is positive for $\varepsilon<1$, because $J_\varepsilon'(0)<0$. So $\psi$ is concave and strictly increasing, with $\varepsilon\le\psi(\varepsilon)\le J_\varepsilon(t)$ for every $t$, whence $\psi(0)=\inf_tD(t)=0$, $\psi(1)=1$ and continuity at the endpoints. Hence $\delta(\cD)=d\psi(\varepsilon)+o(d)$, while $a_p\sqrt d=o(d)$: for every fixed $p$ the condition of Theorem~II holds for large $d$, in the direction of success if $\alpha>\psi(\varepsilon)$ and of failure if $\alpha<\psi(\varepsilon)$.
\end{proof}

\begin{proposition}[Three formulas, one curve]\label[proposition]{prop:lasso:three}
For $\alpha\in(0,1)$ the following coincide: \textup{(a)} $\psi^{-1}(\alpha)$; \textup{(b)} $\alpha\rho_{\mathrm{SE}}(\alpha)$, with the state-evolution formula of \textcite[eq.~(5), signed case]{DonohoMalekiMontanari2009}
\[
\rho_{\mathrm{SE}}(\alpha)=\max_{z\ge0}\frac{1-(2/\alpha)[(1+z^2)\Phi(-z)-z\phi(z)]}{1+z^2-2[(1+z^2)\Phi(-z)-z\phi(z)]};
\]
\textup{(c)} the replica threshold of \textcite[eqs.~(11)--(12)]{KabashimaEtAl2009}: the $\varepsilon$ for which some $t>0$ has $\alpha=J_\varepsilon(t)$ and $\alpha=2(1-\varepsilon)\Phi(-t)+\varepsilon$, with $t=\hat\chi^{-1/2}$ in their notation.
\end{proposition}

\begin{proof}
The plan is to write each formula through $J_\varepsilon$. In (b) the bracket is $D(z)/2$. Since $\psi$ is continuous and increasing and its minimum is attained, $\psi^{-1}(\alpha)$ is the supremum of the $\varepsilon$ with $J_\varepsilon(t)\le\alpha$ for some $t$, and for $t>0$ that condition reads $\varepsilon\le(\alpha-D(t))/(1+t^2-D(t))$, a number below one; exchanging the suprema gives (b). In (c), multiplying equation~(11) of Kabashima et al.\ by $t^2=1/\hat\chi$ turns it into $\alpha=J_\varepsilon(t)$, and $J_\varepsilon'(t)=2\varepsilon t-4(1-\varepsilon)[\phi(t)-t\Phi(-t)]$ gives the identity
\[
J_\varepsilon(t)-\bigl[2(1-\varepsilon)\Phi(-t)+\varepsilon\bigr]=\tfrac t2J_\varepsilon'(t);
\]
so for $t>0$ the two equations say $J_\varepsilon'(t)=0$ and $\alpha=J_\varepsilon(t)$, that is, $\alpha=\psi(\varepsilon)$.
\end{proof}

The formulas appeared separately: \textcite{DonohoMalekiMontanari2009} reported the agreement of state evolution with the polytope curves as a numerical finding; Amelunxen et al.\ note that their upper bound is equivalent to Donoho and Tanner's asymptotic bound and the lower bound is new (their Remark~4.6); \textcite{BayatiLelargeMontanari2015} proved the transition asymptotically by AMP, for independent sub-Gaussian entries.

\section{Measuring it}\label{sec:lasso:measuring}

\begin{figure}[t]
\begin{fullwidth}
\centering
\begin{tikzpicture}
\begin{groupplot}[group style={group size=3 by 1,horizontal sep=15mm},rplot,width=58mm,height=43mm]
\nextgroupplot[xlabel={penalty $\lambda$},ylabel={$\mathrm{MSE}$},xmin=0,xmax=1.05,ymin=0.015,ymax=0.06,
  title={(a) error},ytick={0.02,0.03,0.04,0.05,0.06},yticklabel style={/pgf/number format/precision=2},legend pos=south east]
\addplot[physpred] table[x=lam,y=mse] {figures/data/10-lasso-risk-curve.dat};
\addplot[measured,ebars] table[x=lam,y=mse,y error plus=ci,y error minus=ci] {figures/data/10-lasso-risk-sim.dat};
\legend{fixed point,LASSO}
\nextgroupplot[xlabel={penalty $\lambda$},ylabel={nonzero fraction of $\hat\theta$},xmin=0,xmax=1.05,ymin=0,ymax=0.5,title={(b) sparsity}]
\addplot[guide] coordinates {(0,0.1) (1.05,0.1)};
\addplot[physpred] table[x=lam,y=nnz] {figures/data/10-lasso-risk-curve.dat};
\addplot[measured,ebars] table[x=lam,y=nnz,y error plus=nnzci,y error minus=nnzci] {figures/data/10-lasso-risk-sim.dat};
\node[font=\scriptsize,text=muted,anchor=south west] at (axis cs:0.62,0.1) {true fraction $\varepsilon$};
\nextgroupplot[xlabel={iteration $t$},ylabel={$\mathrm{MSE}$ of AMP},xmin=0,xmax=21,ymin=0.02,ymax=0.05,title={(c) AMP},
  ytick={0.02,0.03,0.04,0.05},legend pos=north east]
\addplot[proofpred] table[x=iter,y=se] {figures/data/10-lasso-amp.dat};
\addplot[measured,ebars] table[x=iter,y=amp,y error plus=ci,y error minus=ci] {figures/data/10-lasso-amp.dat};
\addplot[guide] coordinates {(0,0.023372) (21,0.023372)};
\legend{state evolution,AMP}
\end{groupplot}
\end{tikzpicture}
\end{fullwidth}
\caption{The LASSO for the Bernoulli--Gaussian signal, $\varepsilon=0.1$, $\alpha=0.5$, $\sigma^2=0.05$. (a), (b) Error and fraction of nonzero coordinates against the penalty: carnelian, the fixed point \eqref{eq:lasso:fp}, which \cref{thm:lasso:risk,thm:lasso:cgmt} prove; dots, the LASSO by coordinate descent at $d=1000$, $n=500$, mean of $20$ data sets (seeds $1000$--$1019$, the same for every $\lambda$) with Student-$t$ $95\%$ intervals. (c) AMP with thresholds $t_*\tau_t$ at the optimal $\lambda_*=0.23378$, $d=4000$, mean of $5$ data sets (seeds $3000$--$3004$) with $95\%$ intervals, against state evolution (lapis dashed); gray, the fixed point $0.023372$. Script: \texttt{figures/scripts/10-lasso-risk.py}.}\label{fig:lasso:risk}
\end{figure}

\paragraph{The LASSO.} \texttt{10-lasso-risk.py}\sidenote{Each of the three scripts runs in about a minute and prints every number quoted here. The LASSO is computed by coordinate descent in scikit-learn, whose objective $\frac1{2n}\norm{y-A\theta}^2+a\norm\theta_1$ has $a=\lambda/n$ in our scaling; tolerance $10^{-10}$, warm starts along the grid of $\lambda$.} solves \eqref{eq:lasso:fp} by bisection in $\tau^2$ at each multiplier $t=\zeta/\tau$ above $t_{\min}=0.405234$, the root of $D(t)=\alpha$ \parencite[Proposition~1.3]{BayatiMontanari2012}, and independently computes $\max_b\min_\tau\Psi$; the two agree to eight digits ($\tau_*=0.31103735$, $b_*=0.17155008$). The optimal $\lambda_*=0.23378$ gives $\mathrm{MSE}_*=0.023372$ ($R_*=0.046744$) and keeps $22.4\%$ of the coordinates, twice the $10\%$ that are nonzero. In \cref{fig:lasso:risk}(a,b) all ten means lie within their intervals of the prediction, and at $\lambda_*$ so do the means for $d=250$ to $4000$ ($80$ to $5$ data sets, $0.0229$ to $0.0239$, against $0.02337$). In panel (c) AMP matches state evolution to $10^{-4}$ after ten iterations; its iterates end within $10^{-7}$ (in $\norm\cdot^2/d$) of the LASSO at the data-dependent penalty $\zeta_t(1-\langle\eta'\rangle/\alpha)$ of their fixed point, at most $5.6\%$ from $\lambda_*$, and AMP and the LASSO at $\lambda_*$ give $0.02334\pm0.00157$ and $0.02331\pm0.00157$.

\begin{figure}[t]
\begin{fullwidth}
\centering
\begin{tikzpicture}
\begin{groupplot}[group style={group size=2 by 1,horizontal sep=17mm},rplot,width=82mm,height=43mm]
\nextgroupplot[xlabel={$\delta=\alpha=n/d$},ylabel={$\rho=k/n$},xmin=0,xmax=1,ymin=0,ymax=1,title={(a) the phase diagram},legend pos=north west,legend style={font=\scriptsize}]
\addplot[physpred] table[x=alpha,y=rhoc] {figures/data/10-lasso-transition-curve.dat};
\addplot[proofpred] table[x=alpha,y=rhoc] {figures/data/10-lasso-transition-curve.dat};
\addplot[measured,ebars] table[x=alpha,y=rho50,y error plus=hi,y error minus=lo] {figures/data/10-lasso-transition-pts.dat};
\legend{replica and state evolution,statistical dimension,50\% points ($d=100$)}
\node[font=\scriptsize,text=muted] at (axis cs:0.2,0.08) {recovery};
\node[font=\scriptsize,text=muted] at (axis cs:0.75,0.9) {failure};
\nextgroupplot[xlabel={$\rho=k/n$ at $\alpha=0.5$},ylabel={success fraction},xmin=0.22,xmax=0.58,ymin=-0.05,ymax=1.08,title={(b) finite size},legend pos=north east]
\addplot[guide,forget plot] coordinates {(0.38568967,-0.05) (0.38568967,1.08)};
\addplot[measured3,ebars] table[x=r50,y=p50,y error plus=hi50,y error minus=lo50] {figures/data/10-lasso-transition-prob.dat};
\addplot[measured2,ebars] table[x=r100,y=p100,y error plus=hi100,y error minus=lo100] {figures/data/10-lasso-transition-prob.dat};
\addplot[measured,ebars] table[x=r200,y=p200,y error plus=hi200,y error minus=lo200] {figures/data/10-lasso-transition-prob.dat};
\legend{{$d=50$},{$d=100$},{$d=200$}}
\end{groupplot}
\end{tikzpicture}
\end{fullwidth}
\caption{The Donoho--Tanner transition. (a) The threshold $\rho=\psi^{-1}(\alpha)/\alpha$: the replica and state-evolution formulas (carnelian) and the statistical-dimension curve of \cref{cor:lasso:dt} (lapis dashed) are one function (\cref{prop:lasso:three}). Dots: $50\%$ points of the success probability of $\ell_1$ minimization at $d=100$, from logistic fits to $7$ values of $\rho$ with $24$ problems each (seeds $20000+100j+i$), with $95\%$ parametric-bootstrap intervals ($400$ resamples). (b) Success fraction at $\alpha=0.5$, $30$ problems per point (seeds $10000+1000d/50+i$), Clopper--Pearson $95\%$ intervals; vertical line, $\rho_c(0.5)=0.38568967$. Nonzero entries are standard normal; success means $\norm{\hat\theta-\theta_0}\le10^{-4}\norm{\theta_0}$ for the linear program's solution. Script: \texttt{figures/scripts/10-lasso-transition.py}.}\label{fig:lasso:transition}
\end{figure}

\paragraph{The noiseless threshold.} \texttt{10-lasso-transition.py} evaluates the three formulas of \cref{prop:lasso:three} at $104$ values of $\alpha$; they agree to $1.4\times10^{-14}$. At $\alpha=0.1$, $0.25$, $0.5$, $0.75$, $0.9$ the threshold is $\rho_c=0.18942937$, $0.26738355$, $0.38568967$, $0.53371168$, $0.67816938$, and $\psi(0.5)=0.83129991$, which \textcite{KabashimaEtAl2009} print as $0.83129$. A Monte Carlo estimate of $\delta(\cD)$ at $d=100$ ($20000$ draws) lies between the bounds of \cref{thm:lasso:statdim}, $0.5$ to $0.7$ below $d\psi(s/d)$ for $s=5$ to $50$ (at $s=20$, $50.53\pm0.16$ against $51.11$), where the bound allows $2.8$ to $8.9$. In \cref{fig:lasso:transition} the transition sharpens with $d$, and the $50\%$ points lie on or slightly above the curve, by at most $0.03$ in $\rho$ (two half-widths, at $\alpha=0.85$); with $n=15$ to $85$ measurements this is the expected direction, since $\delta(\cD)$ is below $d\psi(s/d)$.

\begin{figure}[t]
\centering
\begin{tikzpicture}
\begin{axis}[rplot,height=0.37\linewidth,xmode=log,xmin=0.1,xmax=10,ymin=0.35,ymax=1.75,xlabel={Huber threshold $k$},ylabel={$R=\norm{\hat\beta-\beta^*}^2$},
  xtick={0.1,0.3,1,3,10},xticklabels={$0.1$,$0.3$,$1$,$3$,$10$},ytick={0.4,0.6,0.8,1,1.2},legend pos=north west,legend columns=3,
  legend style={font=\scriptsize}]
\addplot[physpred] table[x=k,y=t3] {figures/data/10-lasso-robust-curve.dat};
\addplot[measured,ebars] table[x=k,y=t3,y error plus=tci,y error minus=tci] {figures/data/10-lasso-robust-sim.dat};
\addplot[physpred2] table[x=k,y=gauss] {figures/data/10-lasso-robust-curve.dat};
\addplot[measured2,ebars] table[x=k,y=gauss,y error plus=gci,y error minus=gci] {figures/data/10-lasso-robust-sim.dat};
\addplot[proofpred] coordinates {(0.1,0.642857) (10,0.642857)};
\addplot[guide] coordinates {(0.1,1.285714) (10,1.285714)};
\legend{system ($t_3$),simulation ($t_3$),system (Gaussian),simulation (Gaussian),bound ($t_3$)}
\end{axis}
\end{tikzpicture}
\caption{Huber regression, $p/n=0.3$, noise $t_3$ (variance $3$) or standard Gaussian. Curves: the system \eqref{eq:lasso:mest}. Points: $n=1000$, $p=300$, mean of $10$ data sets per point (seeds $4000+1000i_{\mathrm{noise}}+100i_k+r$), Student-$t$ $95\%$ intervals. Lapis dashed: the information bound $1/((\alpha-1)I)=0.643$ for $t_3$; for Gaussian noise it equals the least-squares risk $3/7$. Gray: least squares under $t_3$, $9/7$. Script: \texttt{figures/scripts/10-lasso-robust.py}.}\label{fig:lasso:robust}
\end{figure}

\paragraph{Robust regression.} \texttt{10-lasso-robust.py} solves \eqref{eq:lasso:mest} for Huber's loss (\cref{sec:lasso:physics}); expectations over $t_3$ noise use a $400$-point Gauss--Legendre rule after $w=\sqrt3\tan u$. Under Gaussian noise $R$ decreases in $k$ to the least-squares value $3/7$, which is also the information bound; the customary $k=1.345$ costs $3.5\%$ ($R=0.44370$). Under $t_3$ noise the best threshold is $k=1.138$, with $R=0.77818$, $39\%$ below least squares ($9/7$) and $21\%$ above the bound $0.64286$. The simulations agree with the system within their intervals at all twelve values of $k$ and for least squares. No measurement contradicts a theorem of the chapter.

\section{The status board}\label{sec:lasso:status}

For Gaussian designs every prediction of the physics box is a theorem, most with two proofs or more; replica symmetry is right because the problems are convex, and the physics rejects its own symmetric solution only for non-convex $\ell_0$ reconstruction.

\begin{statusboard}
\statusrow{LASSO risk and law of $(\hat\theta_j,\theta_{0j})$: the fixed point \eqref{eq:lasso:fp}}{\proved}{Bayati--Montanari 2012, Thm~1.5; Thrampoulidis et al.\ 2018, \S5.4; Miolane--Montanari 2021; \cref{thm:lasso:cgmt}}
\statusrow{Replica-symmetric derivations (LASSO, $\ell_1$, M-estimation)}{\heuristicstatus}{Kabashima et al.\ 2009; Rangan et al.\ 2012; Advani--Ganguli 2016; they agree with the theorems (\cref{prop:lasso:rsao,prop:lasso:three})}
\statusrow{$\ell_1$ weak threshold $\alpha=\psi(\varepsilon)$ (Donoho--Tanner)}{\proved}{Amelunxen et al.\ 2014, Thm~II, Prop.~4.5; Bayati et al.\ 2015}
\statusrow{State evolution, replica and statistical dimension give one curve}{\proved}{\cref{prop:lasso:three}; numerical in Donoho et al.\ 2009}
\statusrow{Replica-symmetric description of $\ell_0$ reconstruction}{\knownfalse}{success solution AT-unstable for every $\alpha$ (Kabashima et al.\ 2009, \S4.1); $\alpha=\varepsilon$ is the counting bound}
\statusrow{LASSO formulas, independent non-Gaussian entries}{\provedunder}{Han--Shen 2023 (sub-Gaussian; large $\lambda$ if $\alpha<1$); $\ell_1$: Bayati et al.\ 2015}
\statusrow{Gaussian formulas for general isotropic designs}{\knownfalse}{Han--Shen 2023: least squares, rows $U_iZ_i$; ridge, LASSO in simulations}
\statusrow{M-estimation system \eqref{eq:lasso:mest}, Gaussian design, $\alpha>1$}{\proved}{Donoho--Montanari 2016, Thm~4.1 (smooth, strongly convex $\rho$); Thrampoulidis et al.\ 2018, Thm~4.1 (convex $\rho$)}
\statusrow{The same, independent non-Gaussian entries}{\provedpreprint}{El Karoui, arXiv:1311.2445v1, Thms~1.1, 6.1 (strongly convex $\rho$); refereed only with a ridge penalty (El Karoui 2018; Han--Shen 2023)}
\statusrow{Variance inflation $R\ge1/((\alpha-1)I(P_W))$}{\proved}{Donoho--Montanari 2016, Cor.~4.3; replicas: Advani--Ganguli 2016}
\end{statusboard}

\section{Where the translation fails}\label{sec:lasso:fails}

\paragraph{What was assumed.} The derivation and every proof here assume independent Gaussian entries: the replica computation when it treats the fields as Gaussian, the CGMT for Gordon's comparison, AMP for its conditioning lemma. \Textcite{KabashimaEtAl2009} argue that their result holds for independent entries of fixed variance; this is a theorem for the $\ell_1$ transition \parencite{BayatiLelargeMontanari2015} and, for sub-Gaussian entries and (if $\alpha<1$) a large enough $\lambda$, for the LASSO \parencite{HanShen2023}. Both sides fix $\lambda,\sigma>0$ and a converging empirical distribution of the signal, so the statements are pointwise in the signal; uniformity over $\ell_p$ balls holds, and over $\ell_0$ balls fails for the empirical law \parencite[Remark~2]{MiolaneMontanari2021}. Our theorem adds Gaussian noise and an independent signal and controls only the norm of the error.

\paragraph{Where the physics is wrong, or not yet right.} For $\ell_0$ reconstruction the replica-symmetric success solution is unstable for every $\alpha$, because hard thresholding is discontinuous \parencite[\S4.1]{KabashimaEtAl2009}; its threshold $\alpha=\varepsilon$ is the counting bound, which exhaustive search attains (\cref{ex:lasso:lzero}), but the symmetric description is wrong, and for non-convex penalties Gordon's comparison gives one inequality only. The fields need not be Gaussian: for isotropic rows $U_iZ_i$, $U_i$ a scalar that is rarely large, \textcite{HanShen2023} prove that the least-squares risk can exceed its Gaussian value by an arbitrary factor, and see the same for ridge and the LASSO in simulations. The thresholds concern a fixed signal; recovering all $k$-sparse vectors at once, the strong threshold, needs more measurements, and a typical-case computation does not address it. Finally, a computation made at $d=\infty$ is silent about the finite-size shifts of \cref{fig:lasso:transition}.

\paragraph{Open problems.} First, the error of the statistical-dimension recipe, which \textcite{AmelunxenEtAl2014} leave open (their Remark~4.4); for $\ell_1$ our Monte Carlo suggests that $\delta(\cD)-d\psi(s/d)$ stays bounded (\cref{ex:lasso:finitesize}). Second, which designs share the Gaussian LASSO risk: independent entries suffice under conditions, general isotropic rows do not, and \textcite{HanShen2023} leave open whether i.i.d.\ rows whose sample covariance obeys the Marchenko--Pastur law suffice. Third, the typical performance of non-convex penalties ($\ell_0$, $\ell_p$ with $p<1$) with noise: the symmetry-breaking computation that \textcite{KabashimaEtAl2009} announced, and a rigorous counterpart. Fourth, a refereed proof of the unregularized robust-regression system for non-Gaussian designs; the refereed results need a ridge penalty.

\begin{exercises}
\item \consolidation (Ridge as a special case.) Replace $\lambda\norm\theta_1$ by $\frac\lambda2\norm\theta^2$. (a) Show that the single site of \cref{comp:lasso:saddle} becomes $s(\Theta_0+\tau Z)$ with $s=1/(1+\lambda/\hat Q)$ and $\hat Q=1-s/\alpha$. (b) Show that $\lambda$ is then exactly the per-sample ridge penalty of \cref{ch:ridge}, and that with $\alpha=2$, $\E\Theta_0^2=2$ (so $r^2=1$), $\sigma^2=0.25$, $\lambda=0.1$ the fixed point and \cref{thm:ridge:risk} both give $R=0.17678$.\label{ex:lasso:ridge}

\item \consolidation (The counting bound.) Let $\theta_0\in\R^d$ be fixed with $k$ nonzero entries and $A\in\R^{n\times d}$ Gaussian with $n\ge k+1$. Show that, almost surely, $\theta_0$ is the only solution of $A\theta=A\theta_0$ with at most $k$ nonzero entries. \hint{for $T\ne\operatorname{supp}\theta_0$ with $\abs T\le k$, the part of $A\theta_0$ carried by columns outside $T$ is independent of the columns in $T$.}\label{ex:lasso:lzero}

\item \consolidation\hard (Variance inflation.) (a) From \eqref{eq:lasso:mest} with $\rho(u)=u^2/2$ derive $R=\sigma^2/(\alpha-1)$. (b) With $\psi_c(z)=z-\operatorname{prox}_c(z)$ and $f_r$ the density of $\hat z$, show $\E\psi_c'(\hat z)=1/\alpha$ and $\E\psi_c(\hat z)^2=r^2/\alpha$, and deduce $r^2\ge1/(\alpha I(f_r))$. (c) With Stam's inequality $1/I(f_r)\ge1/I(P_W)+r^2$ \parencite{Stam1959} conclude $R\ge1/((\alpha-1)I(P_W))$; evaluate it for $t_3$ noise, of Fisher information $\frac23$. \hint{$\E\psi_c'(\hat z)=-\E[\psi_c(\hat z)(\log f_r)'(\hat z)]$, and the Cauchy--Schwarz inequality.}\label{ex:lasso:fisher}

\item \byhand\hard (Replicas for robust regression.) Carry out the computation of \cref{sec:lasso:physics} for $F(\theta)=\sum_i\rho(y_i-\ip{a_i}\theta)$: the energetic term becomes $\alpha\E M_{\chi/\alpha}(W+\sqrt{\cE/\alpha}\,Z)$, the entropic term is that of a Gaussian channel, and the saddle point is \eqref{eq:lasso:mest} with $r^2=\cE/\alpha$, $c=\chi/\alpha$. \hint{$\partial_sM_c(s)=(s-\operatorname{prox}_c(s))/c$, $\partial_cM_c(s)=-(s-\operatorname{prox}_c(s))^2/(2c^2)$, and $\hat Q=\hat m=1/(\alpha c)$ at the saddle.}\label{ex:lasso:robustreplica}

\item \byhand (Very sparse signals.) Show that $\psi(\varepsilon)=2\varepsilon\log(1/\varepsilon)(1+o(1))$ as $\varepsilon\to0$, so that $\rho_c(\delta)\sim1/(2\log(1/\delta))$ as $\delta\to0$, as \textcite{DonohoMalekiMontanari2009} report. The ratio $\psi(\varepsilon)/(2\varepsilon\log(1/\varepsilon))$ is only $0.66$ at $\varepsilon=10^{-2}$ and $0.79$ at $10^{-8}$; explain. \hint{$D(t)\approx4\phi(t)/t^3$, and the minimizer has $t^2=2\log(1/\varepsilon)-O(\log\log(1/\varepsilon))$.}\label{ex:lasso:sparse}

\item \byhand (Nonnegative signals.) For $\norm\cdot_1$ restricted to the nonnegative orthant and $x\ge0$ with $s$ positive entries, show $\delta(\cD)\le d\psi_+(s/d)$, $\psi_+(\varepsilon)=\min_t\{\varepsilon(1+t^2)+(1-\varepsilon)\E(Z-t)_+^2\}$, and check numerically that $\psi_+^{-1}(\alpha)/\alpha$ is the formula of \cref{prop:lasso:three} with $2$ replaced by $1$: $0.240976$, $0.558228$, $0.902657$ at $\alpha=0.1$, $0.5$, $0.9$.\label{ex:lasso:positive}

\item \translate (The overlap concentrates.) Let $\abs{f(v)-f(v')}\le L\norm{v-v'}/\sqrt n$ and $f(\hat v)\to\bar f$ in probability for the minimizer $\hat v$ of $\ell_{b_*}$. Show by steps 4 and 5 of the proof of \cref{thm:lasso:cgmt} that $f(\hat\theta-\theta_0)\to\bar f$, and deduce that the replica overlaps $m=\ip{\hat\theta}{\theta_0}/d$ and $Q=\norm{\hat\theta}^2/d$ tend to $\E[X_*\Theta_0]$ and $\E X_*^2$.\label{ex:lasso:overlap}

\item \translate (The cavity field of AMP.) Show that a fixed point $(x,z)$ of AMP with threshold $\zeta$ satisfies the optimality conditions of the LASSO with $\lambda=\zeta(1-\langle\eta'(x+A^\top z;\zeta)\rangle/\alpha)$, and identify the cavity field with $x+A^\top z$ and $1-\langle\eta'\rangle/\alpha$ with $\hat Q$.\label{ex:lasso:ampfixed}

\item \measurement (Noise sensitivity.) Modify \texttt{10-lasso-risk.py} to minimize the predicted error over $\lambda$ at $\alpha=0.5$ for $\sigma^2=0.05,\dots,0.00005$. For $\varepsilon=0.1<\varepsilon_c(0.5)$, $\mathrm{MSE}_*/\sigma^2$ tends to about $0.96$; for $\varepsilon=0.25$ the error tends to about $0.022$. Explain; compare \textcite{DonohoEtAl2011}.\label{ex:lasso:noisesens}

\item \research (The constant in the statistical dimension.) Estimate $\delta(\cD)-d\psi(s/d)$ by Monte Carlo for $s/d=0.1$ and $d=50$ to $800$. Does it converge, and to what? \textcite[Remark~4.4]{AmelunxenEtAl2014} pose the general question of an optimal error estimate.\label{ex:lasso:finitesize}
\end{exercises}

\notesandrefs
The proportional theory of $\ell_1$ recovery began with polytope geometry: \textcite{Donoho2006} related exact recovery to the neighborliness of random projections of the cross-polytope, and \textcite{DonohoTanner2009a} proved exact thresholds when the projection lowers the dimension radically. \Textcite{Stojnic2009} also obtained the threshold, in an arXiv paper, as a sufficient condition for recovery, from Gordon's escape through a mesh \parencite{Gordon1988} and linear-programming duality; his later combination of Gordon's comparison with convex duality (\cref{ch:gordon}) is the idea that Thrampoulidis, Oymak and Hassibi turned into the CGMT (arXiv 2014; COLT 2015) \parencite{ThrampoulidisOymakHassibi2014,ThrampoulidisOymakHassibi2015}.

The proof of \cref{thm:lasso:cgmt} assembles, in the book's own form, ingredients of \textcite{ThrampoulidisOymakHassibi2015}, \textcite{ThrampoulidisAbbasiHassibi2018} and \textcite{MiolaneMontanari2021}: the absorption of the noise, the function $\Psi$, the strong-convexity gap. The functional $\cF$ of \cref{lem:lasso:scalar} and the change of variables of \cref{prop:lasso:rsao} are the book's own; \cref{prop:lasso:three} collects elementary identities that we do not claim as new; the proof of \cref{thm:lasso:statdim} is that of Amelunxen, Lotz, McCoy and Tropp.

\paragraph{Translating the literature.} Bayati and Montanari write $N$ for our $d$, $\delta$ for $\alpha$, $x_0$ for $\theta_0$, $\alpha$ for our multiplier $t$, and $\theta$ for our threshold $\zeta$. Donoho, Maleki and Montanari use the same $\delta$ and write $\rho=k/n$, our $\varepsilon/\alpha$. Kabashima, Wadayama and Tanaka write $P$, $N$ for our $n$, $d$, $\rho$ for $\varepsilon$, and $F$, with entries of variance $1/N$, for $A$; their conjugate fields are ours, and their $\phi_p(h;\hat Q)$ is our single site with $\lambda=1$. Amelunxen et al.\ write $m$ for our $n$, $s$ for $k$, $\tau$ for $t$. Miolane and Montanari write $\beta$ for our $b$ and $s_{\max}(\delta)$ for $\varepsilon_c(\alpha)$. El Karoui writes $\kappa=p/n$, $r_\rho$, $c_\rho$, and Donoho and Montanari $\delta=n/p$, $\tau$, $b$, for our $1/\alpha$, $r$, $c$. Theorem and section numbers of \textcite{ThrampoulidisAbbasiHassibi2018} are those of its arXiv version~1.

\printchapterbib
\end{refsection}
